% ID: SEL-D0004-OA-E03
% ESTADO: cotejada. Cotejo literal independiente y revisión de composición realizados.
% Fuente: Selectividad/sel_2019_matematicas_aplicadas/Suplente Septiembre/Exámenes Matemáticas Aplicadas CCSS - A y B.pdf
% SHA-256: eecad753694b3d8cb2d427efddb00d75faae9c4ec5b027056291bc95d7d2e3da
% Página: 1; opción A; ejercicio 3.
% Encabezado original: EJERCICIO 3. Solo enunciado; no contiene solución.

En una localidad andaluza hay tres institutos de ESO. De los 500 estudiantes
que cursan 1\textsuperscript{o} de ESO en dicha localidad, 250 están
matriculados en el instituto A, 150 en el B y el resto están matriculados
en el instituto C. Se sabe que han superado la materia de Matemáticas el
70\,\% del alumnado de 1\textsuperscript{o} de ESO matriculado en el
instituto A, el 68\,\% de B y el 73\,\% de C. Se elige al azar un
estudiante de 1\textsuperscript{o} de ESO de la citada localidad.

\begin{enumerate}[label=\alph*)]
  \item \textbf{(1 punto)} Calcule la probabilidad de que no haya superado
  Matemáticas.
  \item \textbf{(1 punto)} Calcule la probabilidad de que esté matriculado
  en el instituto A, sabiendo que ha superado Matemáticas.
  \item \textbf{(0.5 puntos)} Calcule la probabilidad de que esté
  matriculado en el instituto C y no haya superado Matemáticas.
\end{enumerate}
